% !TeX root = ../example.tex
% Keep these examples within the existing subsection/navigation structure.
\begin{frame}[fragile]{代码展示：Python}
\begin{campuscode}[language=Python,numbers,caption={保留缩进，便于讲解函数结构。}]{normalize.py}
def normalize(values):
    # 将数值归一化，保留清晰的缩进
    total = sum(values)
    if total == 0:
        return [0.0 for value in values]
    return [value / total for value in values]

print(normalize([2, 3, 5]))
\end{campuscode}
\small 标题栏标识文件与语言；下方说明补充阅读提示。
\end{frame}

\begin{frame}[fragile]{代码展示：终端命令与输出}
\begin{campuscode}[language=bash,caption={静态展示命令与输出，不在编译时执行。}]{Terminal}
$ python -m venv .venv
$ source .venv/bin/activate
$ python normalize.py
[0.2, 0.3, 0.5]
\end{campuscode}
\small 用命令提示符区分输入与输出。
\end{frame}

\begin{frame}[fragile]{代码展示：从文件读取}
\campusinputcode[language=Python,description={从文件读取并高亮选定源码。}]
  {normalize.py}{chapters/code/normalize.py}
\small 用 \verb|\campusinputcode[language=Python]{文件名}{文件路径}| 复用源文件。
\end{frame}
